\documentclass[10pt,a4paper]{amsart}
\usepackage{amsmath,amssymb,amsthm,graphicx}
\usepackage{verbatim}
\newenvironment{reference}{\comment}{\endcomment}
%\newenvironment{reference}{}{}
\newenvironment{slogan}{\comment}{\endcomment}
\newenvironment{history}{\comment}{\endcomment}

% For commutative diagrams we use Xy-pic
\usepackage[all]{xy}

% We use 2cell for 2-commutative diagrams.
\xyoption{2cell}
% Theorem environments.
%
\theoremstyle{plain}
\newtheorem{theorem}[subsection]{Theorem}
\newtheorem{proposition}[subsection]{Proposition}
\newtheorem{lemma}[subsection]{Lemma}

\theoremstyle{definition}
\newtheorem{definition}[subsection]{Definition}
\newtheorem{example}[subsection]{Example}
\newtheorem{exercise}[subsection]{Exercise}
\newtheorem{situation}[subsection]{Situation}

\theoremstyle{remark}
\newtheorem{remark}[subsection]{Remark}
\newtheorem{remarks}[subsection]{Remarks}

\numberwithin{equation}{subsection}

% Macros
%
\def\lim{\mathop{\mathrm{lim}}\nolimits}
\def\colim{\mathop{\mathrm{colim}}\nolimits}
\def\Spec{\mathop{\mathrm{Spec}}}
\def\Hom{\mathop{\mathrm{Hom}}\nolimits}
\def\Ext{\mathop{\mathrm{Ext}}\nolimits}
\def\SheafHom{\mathop{\mathcal{H}\!\mathit{om}}\nolimits}
\def\SheafExt{\mathop{\mathcal{E}\!\mathit{xt}}\nolimits}
\def\Sch{\mathit{Sch}}
\def\Mor{\mathop{\mathrm{Mor}}\nolimits}
\def\Ob{\mathop{\mathrm{Ob}}\nolimits}
\def\Sh{\mathop{\mathit{Sh}}\nolimits}
\def\NL{\mathop{N\!L}\nolimits}
\def\CH{\mathop{\mathrm{CH}}\nolimits}
\def\proetale{{pro\text{-}\acute{e}tale}}
\def\etale{{\acute{e}tale}}
\def\QCoh{\mathit{QCoh}}
\def\Ker{\mathop{\mathrm{Ker}}}
\def\Im{\mathop{\mathrm{Im}}}
\def\Coker{\mathop{\mathrm{Coker}}}
\def\Coim{\mathop{\mathrm{Coim}}}

% Boxtimes
%
\DeclareMathSymbol{\boxtimes}{\mathbin}{AMSa}{"02}

%
% Macros for moduli stacks/spaces
%
\def\QCohstack{\mathcal{QC}\!\mathit{oh}}
\def\Cohstack{\mathcal{C}\!\mathit{oh}}
\def\Spacesstack{\mathcal{S}\!\mathit{paces}}
\def\Quotfunctor{\mathrm{Quot}}
\def\Hilbfunctor{\mathrm{Hilb}}
\def\Curvesstack{\mathcal{C}\!\mathit{urves}}
\def\Polarizedstack{\mathcal{P}\!\mathit{olarized}}
\def\Complexesstack{\mathcal{C}\!\mathit{omplexes}}
% \Pic is the operator that assigns to X its picard group, usage \Pic(X)
% \Picardstack_{X/B} denotes the Picard stack of X over B
% \Picardfunctor_{X/B} denotes the Picard functor of X over B
\def\Pic{\mathop{\mathrm{Pic}}\nolimits}
\def\Picardstack{\mathcal{P}\!\mathit{ic}}
\def\Picardfunctor{\mathrm{Pic}}
\def\Deformationcategory{\mathcal{D}\!\mathit{ef}}
\usepackage{cleveref}
\crefname{originallemma}{original lemma}{original lemmas}
\Crefname{originallemma}{Original lemma}{Original lemmas}
\makeatletter
\newlabel{lemma-functoriality-differentials}{{08TR}{source}{}{sourcepointer-1}{}}
\newlabel{lemma-functoriality-differentials@cref}{{[originallemma][0][]08TR}{[1][0][]source}{}{}{}}
\makeatother
\crefname{originalreference}{original reference}{original references}
\Crefname{originalreference}{Original reference}{Original references}
\makeatletter
\newlabel{algebra-lemma-differential-seq}{{00RU}{source}{}{sourcepointer-2}{}}
\newlabel{algebra-lemma-differential-seq@cref}{{[originalreference][0][]00RU}{[1][0][]source}{}{}{}}
\makeatother
\crefname{originallemma}{original lemma}{original lemmas}
\Crefname{originallemma}{Original lemma}{Original lemmas}
\makeatletter
\newlabel{lemma-differentials-sheafify}{{08TP}{source}{}{sourcepointer-3}{}}
\newlabel{lemma-differentials-sheafify@cref}{{[originallemma][0][]08TP}{[1][0][]source}{}{}{}}
\makeatother
\title{The conormal sequence \$\textbackslash{}mathcal\{I\}/\textbackslash{}mathcal\{I\}\textasciicircum{}2 \textbackslash{}to \textbackslash{}Omega\_\{\textbackslash{}mathcal\{O\}\_2/\textbackslash{}mathcal\{O\}\_1\} \textbackslash{}otimes \textbackslash{}mathcal\{O\}'\_2 \textbackslash{}to \textbackslash{}Omega\_\{\textbackslash{}mathcal\{O\}'\_2/\textbackslash{}mathcal\{O\}\_1\} \textbackslash{}to 0\$}
\author{The Stacks Project Authors}
\date{stacks-project:sites-modules.tex:local-snapshot-20261012; GNU FDL 1.2 or later}
\begin{document}
\maketitle
\noindent\textit{Editorial note.} All delivered mathematical text is the author's own, copied verbatim from the pinned Stacks Project chapter file sites-modules.tex; no statement or proof step was rewritten, reordered or removed.\par\medskip
\noindent\textit{Editorial note.} Span recomputation by label pairing (LANE\_BRIEF 10.2): in this chapter the author's proof environment follows the statement environment as a standalone block; the spec covers the complete environment pair (statement span, proof span) as recorded in the pool candidate and re-verified against the pinned source (sm-twins-20261014.json span sanity: all OK).\par\medskip
\noindent\textit{Editorial note.} Occupation scan: combined claim\_key + statement\_anchor + generic-line scan over all registration files (active and withdrawn) found hard occupation only from the 47 pre-existing lane-s11 S11-sites-mod-* cards (same pinned snapshot, delivered) and cross-chapter citation noise, 0 other-lane hits for sites-modules.tex (sm-occupied.json, 2026-10-14); the RESUME-pool dispatch unlocks the remaining face.\par\medskip
\ifdefined\hypertarget\hypertarget{sourcepointer-1}{}\fi
\noindent\textit{Source reference.} Original lemma 08TR: sources/stacks-chow/tags.txt, source lines 11174--11174. Cross-chapter standard prerequisite cited by the author inside the counted span; it is used as a pointer and its own proof is not part of this item.\par\medskip
\ifdefined\hypertarget\hypertarget{sourcepointer-2}{}\fi
\noindent\textit{Source reference.} Original reference 00RU: sources/stacks-chow/tags.txt, source lines 1016--1016. Cross-chapter standard prerequisite cited by the author inside the counted span; it is used as a pointer and its own proof is not part of this item.\par\medskip
\ifdefined\hypertarget\hypertarget{sourcepointer-3}{}\fi
\noindent\textit{Source reference.} Original lemma 08TP: sources/stacks-chow/tags.txt, source lines 11172--11172. Cross-chapter standard prerequisite cited by the author inside the counted span; it is used as a pointer and its own proof is not part of this item.\par\medskip
\begin{lemma}
\label{lemma-differential-seq}
In Lemma \ref{lemma-functoriality-differentials} suppose that
$\mathcal{O}_2 \to \mathcal{O}'_2$ is surjective with kernel
$\mathcal{I} \subset \mathcal{O}_2$ and assume that
$\mathcal{O}_1 = \mathcal{O}'_1$. Then there is a canonical exact
sequence of $\mathcal{O}'_2$-modules
$$
\mathcal{I}/\mathcal{I}^2
\longrightarrow
\Omega_{\mathcal{O}_2/\mathcal{O}_1} \otimes_{\mathcal{O}_2} \mathcal{O}'_2
\longrightarrow
\Omega_{\mathcal{O}'_2/\mathcal{O}_1}
\longrightarrow
0
$$
The leftmost map is characterized by the rule that a local section
$f$ of $\mathcal{I}$ maps to $\text{d}f \otimes 1$.
\end{lemma}

\begin{proof}
For a local section $f$ of $\mathcal{I}$ denote $\overline{f}$ the image of
$f$ in $\mathcal{I}/\mathcal{I}^2$. To show that the map
$\overline{f} \mapsto \text{d}f \otimes 1$ is well defined we just have to
check that $\text{d} f_1f_2 \otimes 1 = 0$ if $f_1, f_2$ are local sections
of $\mathcal{I}$. And this is clear from the Leibniz rule
$\text{d} f_1f_2 \otimes 1 =
(f_1 \text{d}f_2 + f_2 \text{d} f_1 )\otimes 1 =
\text{d}f_2 \otimes f_1 + \text{d}f_2 \otimes f_1 = 0$.
A similar computation show this map is
$\mathcal{O}'_2 = \mathcal{O}_2/\mathcal{I}$-linear. The map on the right
is the one from Lemma \ref{lemma-functoriality-differentials}.

\medskip\noindent
To see that the sequence is exact, we argue as follows. Let
$\mathcal{O}''_2 \subset \mathcal{O}'_2$ be the presheaf of
$\mathcal{O}_1$-algebras whose value on $U$ is the image of
$\mathcal{O}_2(U) \to \mathcal{O}'_2(U)$. By
Algebra, Lemma \ref{algebra-lemma-differential-seq}
the sequences
$$
\mathcal{I}(U)/\mathcal{I}(U)^2
\longrightarrow
\Omega_{\mathcal{O}_2(U)/\mathcal{O}_1(U)}
\otimes_{\mathcal{O}_2(U)} \mathcal{O}''_2(U)
\longrightarrow
\Omega_{\mathcal{O}''_2(U)/\mathcal{O}_1(U)}
\longrightarrow
0
$$
are exact for all objects $U$ of $\mathcal{C}$. Since sheafification is
exact this gives an exact sequence of sheaves of
$(\mathcal{O}'_2)^\#$-modules.
By Lemma \ref{lemma-differentials-sheafify}
and the fact that $(\mathcal{O}''_2)^\# = \mathcal{O}'_2$
we conclude.
\end{proof}

\end{document}
